% ===============================================================
%  Chapter 1 -- Introduction and Structure of the Report
% ===============================================================
\chapter{Introduction and Structure of the Report}
\label{ch:intro}

\begin{figure}[htbp]
\centering
\adjustbox{max width=\linewidth}{%
\begin{forest}
  chapter map,
  [Chapter 1
    [Motivation
      [Private caches hold copies]
      [AMD and Intel chose differently]]
    [Scope
      [Write-invalidate family]
      [2 to 16 cores]]
    [Organisation
      [Seven Parts]]
    [Textbook mapping
      [Sarangi Ch.\ 11--13]]
  ]
\end{forest}}
\caption[]{Structure of Chapter~1. Each coloured branch
is one section; the pale boxes are the points it makes.}
\label{fig:ch1-map}
\end{figure}

\section{Motivation}
\label{sec:intro-motivation}
\sectionmap{[Motivation [Every chip is a multicore] [Each core has a private cache]
  [Copies must agree] [Vendors chose differently]]}
\index{MOESI}%

Phones, laptops and servers now all use \emph{multicore}\index{multicore processor}
processors: several cores on one chip, running parts of one program and sharing
one memory. Each core has a small private cache close to it, because one cache
serving every core would be too slow. Once two of these caches hold copies of the
same data, the hardware has to stop the copies from disagreeing in any way a
program could notice. The mechanism that does this is a
\emph{\gls{coherence} protocol}\index{cache coherence protocol}.

The two largest x86 vendors extended the same four-state protocol in different
directions. AMD documents a protocol called \gls{moesi}~\cite{amd64vol2}, and
Intel's QuickPath processors use one called \gls{mesif}~\cite{goodman2009mesif}.
This report looks at the AMD side of that choice: what MOESI's extra state saves
compared with \gls{mesi}, and under which workloads the saving is large enough to
matter.

\section{Scope of the Report}
\label{sec:intro-scope}
\sectionmap{[Scope [In: snooping {\&} directory MSI family] [In: 2--16 cores]
  [Out: memory consistency] [Out: message passing]]}
\index{MESI}\index{MOESI}\index{MSI}\index{snooping protocol}\index{directory protocol}\index{memory consistency}\index{verification}%

The report covers hardware cache coherence in shared-memory multicore processors,
limited to the \gls{invalidate} family that contains \gls{msi}, MESI and MOESI. It
explains the architecture needed to follow these protocols, the inefficiency that
MOESI was designed to remove, the published work on that inefficiency, and our own
simulation study on 2, 4, 8 and 16 cores.

We leave out three neighbouring topics except where a boundary needs marking.
Memory consistency models deal with ordering across different addresses rather
than a single address~\cite[\S12.4.3]{sarangi2025basic}. Message-passing systems
have no shared memory to keep coherent. Formal verification of protocols needs
tools beyond this course.

\section{Organisation of the Report}
\label{sec:intro-parts}
\sectionmap{[Seven Parts [I builds the base] [II states the problem]
  [III reviews prior work] [IV adds our idea] [V--VII test and conclude]]}
\index{MESI}\index{MOESI}%

The report has seven Parts, and each one depends only on the Parts before it
(\cref{fig:report-map}). MESI and MOESI are first described in Part~I
as plain pieces of architecture. We judge them as solutions only in Part~III,
after Part~II has stated the problem they are judged against.

\begin{figure}[htbp]
\centering
\adjustbox{max width=\linewidth}{%
\begin{forest}
  part map,
  [This report
    [Part I: Foundations
      [Caches] [Parallel processing] [Coherence {\&} the MSI family]]
    [Part II: The Problem
      [Origin of overhead] [Forced writebacks]]
    [Part III: Existing Solutions
      [Protocols {\&} papers] [Comparison]]
    [Part IV: Our Contribution
      [Controlled method] [DSR predictor]]
    [Part V: Simulation
      [Two-layer simulator] [Workloads]]
    [Part VI: Results
      [Predictions H1--H4 tested]]
    [Part VII: Limitations {\&} Conclusion
      [Threats] [Answers to RQ1--RQ4]]
  ]
\end{forest}}
\caption[]{The whole report as one map. Each Part
opens with its own map (blue root), each chapter with a chapter map (slate root)
and each section with a small section map (orange root). Decision trees read top
to bottom; in protocol diagrams blue marks MESI and orange marks MOESI.}
\label{fig:report-map}
\end{figure}

\section{Mapping to the Prescribed Textbook}
\label{sec:intro-textbook}
\sectionmap{[Textbook mapping [Chapter 11: memory system] [Chapter 12: multiprocessors]
  [Gap: MESI is an exercise] [Gap: MOESI absent]]}
\index{MESI}\index{MOESI}\index{MSI}\index{Exclusive state}\index{write-update}\index{Sarangi textbook}%

The prescribed textbook is Sarangi's \emph{Basic Computer Architecture},
version~3.09~\cite{sarangi2025basic}. \Cref{tab:textbook-map} lists the textbook
sections each chapter draws on. The textbook develops write-update and MSI
coherence in detail and introduces the Exclusive state in one paragraph, but it
leaves the full MESI protocol ``as an exercise for the
reader''~\cite[\S12.4.6, p.~603]{sarangi2025basic} and does not cover MOESI.
Our report starts from that point.

\begin{table}[htbp]
\centering
\caption{Chapters of this report mapped to sections and pages of the prescribed
textbook~\cite{sarangi2025basic}. ``---'' marks chapters that build on earlier
chapters rather than on new textbook material.}
\label{tab:textbook-map}
\small
\begin{tabular}{@{}L{4.0cm} L{6.6cm} L{2.2cm}@{}}
\toprule
Report chapter & Textbook sections used & Pages \\
\midrule
2 Caches & \S7.4.1--7.4.2 SRAM, CAM; \S11.1--11.2 locality, caches; \S11.3.1--11.3.2 AMAT, misses & 286--537 \\
3 Parallel processing & \S12.1--12.3 background, Amdahl, Flynn; \S12.4.4; \S12.6 interconnects & 567--630 \\
4 Coherence & \S12.4.1--12.4.2 axioms; \S12.4.6 snoopy, directory; \S13.4.3 arbitration & 579--603, 674--676 \\
5 MSI family & \S12.4.6 write-invalidate MSI, Exclusive state & 600--603 \\
6 History & \S12.1 history of multiprocessors & 567--570 \\
7 Origin of the cost & \S12.4.2 coherence; \S12.4.6 bus messages & 581--602 \\
8 Scaling & \S11.3 memory model; \S12.2.3 Amdahl; \S12.6.2 bisection bandwidth & 533--623 \\
9 Dirty sharing & \S11.2.6 evict; \S12.4.6 Fig.~12.14 & 530--602 \\
10 Problem statement & builds on Chapters 7--9 & --- \\
11--12 Owner, update, MESIF & \S12.4.6 write-update, write-invalidate, Exclusive state & 597--603 \\
13 Directories, tokens & \S12.4.6 directories; \S12.6.2--12.6.3 topologies & 599--603, 622--629 \\
14 Comparison & builds on 11--13; \S13.4.3 arbitration & 674--676 \\
15--16 Our method, DSR & \S12.4.1--12.4.2 axioms; \S12.4.6 Fig.~12.14; \S11.3.1 AMAT & 533--534, 579--603 \\
17--18 Simulator, workloads & \S11.2, \S11.2.6, \S11.3.1; \S12.2; \S13.4.3 & 517--534, 570--577, 674--676 \\
19--20 Results & \S11.3.1 AMAT; \S12.6.2 bandwidth; \S12.4.6 Fig.~12.14 & 517--534, 602--622 \\
21--22 Limitations, conclusion & \S11.1--11.3; \S12.4--12.6 & 506--538, 579--629 \\
\bottomrule
\end{tabular}
\end{table}
